\documentclass[11pt,a4paper]{article}
\usepackage[T1]{fontenc}
\usepackage{lmodern}
\usepackage[margin=26mm,headheight=15pt]{geometry}
\usepackage{amsmath,amssymb,amsthm,mathtools,bm,mathrsfs}
\usepackage{microtype,booktabs,longtable,array,enumitem,xcolor,xurl,fancyhdr}
\usepackage[colorlinks=true,linkcolor=blue!40!black,citecolor=blue!40!black,urlcolor=blue!40!black]{hyperref}
\usepackage{bookmark}
\setlength{\emergencystretch}{3em}
\setcounter{tocdepth}{2}
\numberwithin{equation}{section}
\newtheorem{theorem}{Theorem}[section]
\newtheorem{proposition}[theorem]{Proposition}
\newtheorem{lemma}[theorem]{Lemma}
\newtheorem{corollary}[theorem]{Corollary}
\theoremstyle{definition}
\newtheorem{definition}[theorem]{Definition}
\newtheorem{example}[theorem]{Example}
\newtheorem{remark}[theorem]{Remark}
\DeclareMathOperator{\Li}{Li}
\DeclareMathOperator{\Res}{Res}
\DeclareMathOperator{\Rea}{Re}
\DeclareMathOperator{\CT}{CT}
\newcommand{\dd}{\,\mathrm d}
\newcommand{\ii}{\mathrm i}
\newcommand{\C}{\mathbb C}
\newcommand{\Q}{\mathbb Q}
\newcommand{\N}{\mathbb N}
\newcommand{\Z}{\mathbb Z}
\newcommand{\bmhC}{\mathcal C}
\newcommand{\bmhK}{\mathcal K}
\newcommand{\bmhH}{\mathcal H}
\newcommand{\bmhM}{\mathcal M}
\newcommand{\src}[1]{{\small\nolinkurl{#1}}}
\newcommand{\zetad}{\zeta}
\newcommand{\sh}{\mathbin{\amalg}}
\pagestyle{fancy}
\fancyhf{}
\fancyhead[L]{\small Bilinear Mellin--Hurwitz calculus}
\fancyhead[R]{\small ProveIt research continuation}
\fancyfoot[C]{\thepage}
\renewcommand{\headrulewidth}{0.3pt}
\title{\bfseries Bilinear Mellin--Hurwitz Calculus\\[5pt]
\large Double-zeta reductions, harmonic moments,\\and Stieltjes reciprocity}
\author{Research continuation prepared for the ProveIt project\\
\small Mathematical derivations and reproducible checks with OpenAI assistance}
\date{10 October 2026}
\hypersetup{pdftitle={Bilinear Mellin-Hurwitz Calculus},pdfauthor={Research continuation for the ProveIt project}}
\begin{document}
\maketitle
\thispagestyle{empty}
\begin{abstract}
We develop an exact calculus for products of polylogarithms under rational
Mellin kernels. For positive integer orders $m,n$, the transform
$\int_0^\infty x^{a-1}\Li_m(-x)\Li_n(-x)/(1+x)\dd x$ is reduced to a finite
combination of common-shift double Hurwitz zeta functions. Connected
combinations remove the divergent leading-index-one coordinates and make
both resonances $a=0,-1$ explicit. At $a=0$ we obtain an all-weight formula
using only ordinary and double zeta values, together with every logarithmic
moment. A terminating Laurent-kernel recurrence evaluates all admissible
integer Mellin parameters and denominator orders; its elementary base
cases are Gamma derivatives and Bell polynomials in generalized harmonic
numbers. A separate inversion argument reduces the odd-parity central
moments to ordinary zeta data and gives bilinear identities for every
polylogarithm order derivative at one in terms of generalized Stieltjes
constants. We also give explicit iterated-integral primitives, a weighted
multiple-zeta reduction, and a constructive multiscale closure theorem.
The latter settles the positive-integer-order finite-closure part of a
specific research question in the repository; arbitrary spectral orders
and arithmetic minimal depth remain separate questions. Full analytic
proofs, finite exact verification, numerical diagnostics, and integration
notes are supplied. Historical priority beyond the examined sources is
not asserted.
\end{abstract}
\vfill
\noindent\small\textbf{Source scope.} The continuation targets the Mellin--Lerch
report and the canonical Polylogarithms manuscript. The source pin is
\src{b1df802799851c62a860570c3ac78ee31b9f5c04}; the precise inspected files,
unchanged research-question blob, and limitations of the incoming-archive
inspection are recorded in Section~\ref{bmh:sec:scope} and the delivery notes.
\clearpage
\tableofcontents
\clearpage
\section{The problem, conventions, and scope}\label{bmh:sec:scope}

\subsection{A precise continuation of the repository}
The canonical manuscript \emph{Polylogarithms and their Arithmetic Bridges}
collects an extensive theory of polylogarithmic identities, harmonic sums,
Gamma arithmetic, Hurwitz jets, and Stieltjes correlations~\cite{Repo}.
The related report \emph{Mellin--Lerch Transforms and Dilation Identities for
Stieltjes Correlations} proves a one-factor Mellin transform and then asks,
in its tenth proposed research question, for a treatment of
\begin{equation}\label{bmh:eq:problem}
 \int_0^\infty\frac{x^{a-1}\Li_s(-zx)\Li_t(-wx)}{(1+x)^N}\dd x.
\end{equation}
Its two explicit objectives are finite multiple-polylogarithm closure at
integer Mellin parameters and a determination of the depth forced by
independent arguments~\cite{MLD}. The present article answers the first
objective when the polylogarithm orders are positive integers. It goes
further at equal arguments by deriving a finite double-Hurwitz formula
for noninteger $a$, exact formulas at both integer resonances in the
$N=1$ strip, and all logarithmic moments.

This is not a claim to settle the question for arbitrary complex spectral
orders, nor to prove numerical or motivic minimal depth. Generic
hyperlogarithmic integration is established mathematics; the multiscale
closure proof below is an explicit specialization of that method, with
its alphabet, full product convergence strip, and residue-weight
refinement made concrete~\cite{Panzer}. The connected double-Hurwitz
formula and the coefficient identities are derived here in detail and
are the main research contribution relative to the inspected repository
material. An exhaustive literature priority search has not been made.

The main source path is
\begin{center}\small
\url{Analysis/Polylogarithms/docs/reports/stieltjes-correlation-calculus/continuations/mellin-lerch-dilation/}.
\end{center}
The research-question file is \src{sections/06-audit-research.tex}; its
blob at the recorded pin is
\src{c0e3d58e42005ca844ca8b6900d7feaf1b33e2f8}.
The existing one-factor theorem, used with attribution in
Sections~\ref{bmh:sec:reciprocity}--\ref{bmh:sec:stieltjes}, is in
\src{sections/02-mellin.tex}. We also inspected the canonical README,
main source, selected depth and discovery sections, the incoming
inventory, and its intake instructions. The binary contents of the eight
archives in that inventory were not analytically audited. Consequently
our comparison is with the specified inspected text, not a claim of
exhaustive novelty against every incoming archive.

\subsection{Notation and normalization}
All unqualified $m,n,N$ are positive integers; $h,r$ are nonnegative
integers. Put $W=m+n$ and define
\begin{equation}\label{bmh:eq:M-definition}
 M_N(a;m,n;z,w)=\int_0^\infty
 \frac{x^{a-1}\Li_m(-zx)\Li_n(-wx)}{(1+x)^N}\dd x.
\end{equation}
We suppress $z,w$ when both are one. Powers of positive real numbers use
the real logarithm. The polylogarithm is on its principal branch; positive
$z,w$ therefore never put its argument on the positive-real branch cut.
More general scales are considered on
$\C\setminus(-\infty,0]$. Standard polylogarithm, Hurwitz-zeta, and Gamma
conventions are those of the DLMF~\cite{DLMFpoly,DLMFhurwitz,DLMFbeta}.

We use a semicolon in $\zeta(s;q)$ for the single Hurwitz parameter,
to distinguish it from the comma-separated multiple-zeta indices.
Multiple zeta values use \emph{decreasing} indices:
\begin{equation}\label{bmh:eq:MZV}
 \zeta(k_1,\ldots,k_d)=
 \sum_{n_1>\cdots>n_d\ge1}\frac1{n_1^{k_1}\cdots n_d^{k_d}},
 \qquad k_1\ge2,
\end{equation}
and the common-shift double Hurwitz value is
\begin{equation}\label{bmh:eq:Zq}
 Z_q(A,B)=\sum_{\nu>\mu\ge0}
 \frac1{(\nu+q)^A(\mu+q)^B},
 \qquad A\ge2,
 \quad B\ge1,\quad \Rea q>0.
\end{equation}
Thus $Z_1(A,B)=\zeta(A,B)$. When the spectral indices vary in a
proof, the same series is used wherever it is absolutely convergent. Each of these series, and any fixed number of
its $q$-derivatives, converges normally on compact subsets of
$\Rea q>0$. The inner sum is at most logarithmic when $B=1$ and bounded
when $B>1$; this proves the assertion directly. Spectral differentiation
in $B$ near one adds only logarithms and is also normally convergent.

We write $(s)_j=s(s+1)\cdots(s+j-1)$, with $(s)_0=1$, and
$\binom uv=0$ for integer $u\ge0$ and $v\notin\{0,\ldots,u\}$.
Generalized harmonic numbers are
$H_d^{(r)}=\sum_{j=1}^d j^{-r}$, with $H_0^{(r)}=0$.
The Stieltjes convention is
\begin{equation}\label{bmh:eq:Stieltjes-convention}
 \zeta(s;q)=\frac1{s-1}
 +\sum_{r=0}^\infty\frac{(-1)^r}{r!}\gamma_r(q)(s-1)^r,
 \qquad \gamma_0(q)=-\psi(q).
\end{equation}
An order derivative $\Li_s^{[r]}$ differentiates $s$, not the argument.

\subsection{What the principal results say}
Theorem~\ref{bmh:thm:master} is a finite formula at all positive integer
orders and noninteger Mellin parameters in the full product strip.
Theorem~\ref{bmh:thm:compact} gives an especially short unshifted double-zeta
answer. Theorems~\ref{bmh:thm:jets} and \ref{bmh:thm:kernel} give the complete
logarithmic and integer-kernel calculi. Theorem~\ref{bmh:thm:central}
identifies a parity sector with a single-zeta answer, and
Theorem~\ref{bmh:thm:Stieltjes} connects bilinear polylogarithmic integrals
to every generalized Stieltjes index. The algebraic checks are finite
implementation tests. The proofs of the infinite families are the
analytic and combinatorial arguments printed below.

\section{Convergence and constructive multiscale closure}\label{bmh:sec:closure}

\subsection{The product strip}
\begin{proposition}[Full common convergence strip]\label{bmh:prop:strip}
For $r$ polylogarithm factors with scales in
$\C\setminus(-\infty,0]$, the ordinary integral
\[
 \bmhM_N(a;\boldsymbol s;\boldsymbol z)=
 \int_0^\infty\frac{x^{a-1}}{(1+x)^N}
       \prod_{j=1}^r\Li_{s_j}(-z_jx)\dd x
\]
is holomorphic for $-r<\Rea a<N$ and entire in each spectral order.
Arbitrary fixed spectral derivatives and powers of $\log x$ can be
inserted under the integral on compact subsets of this domain.
This is the maximal strip common to all spectral orders.
\end{proposition}
\begin{proof}
At zero the power series gives $\Li_s(-zx)=O(x)$ uniformly on compact
sets of $s,z$. At infinity the bound is $O((1+\log x)^L)$ for some
compact-dependent $L$. Here is a direct justification of the latter
bound that includes nonpositive spectral orders. Choose an integer $k$
such that $\Rea(s+k)>0$ on the spectral compact set. The Fermi integral
\[
 \Li_{s+k}(-zx)=-\frac1{\Gamma(s+k)}\int_0^\infty
 u^{s+k-1}\frac{zx e^{-u}}{1+zx e^{-u}}\dd u
\]
can be differentiated $k$ times by $x\partial_x$. The differentiated
rational kernels are bounded by a constant times
$\min(xe^{-u},1)$ on the chosen scale compact. Split the $u$ integral at
$\log x$. This proves a fixed power-logarithmic bound. Small circles in
the spectral variables and Cauchy's formula give the same conclusion
for fixed spectral derivatives. The resulting majorants are integrable
on every smaller strip; parameter holomorphy and differentiation follow.
Finally $s_j=0,z_j=1$ makes the product $(-x/(1+x))^r$, forcing both
strict endpoint inequalities.
\end{proof}

\begin{proposition}[Denominator recurrence]\label{bmh:prop:Nrec}
With $E_j$ lowering the $j$th spectral order by one,
\begin{equation}\label{bmh:eq:Nrec}
 \bmhM_{N+1}(a)=\frac1N\left(N-a-\sum_{j=1}^rE_j\right)\bmhM_N(a)
\end{equation}
on the common strip $-r<\Rea a<N$.
\end{proposition}
\begin{proof}
Integrate the derivative of
$x^a(1+x)^{-N}\prod_j\Li_{s_j}(-z_jx)$. Its endpoint values vanish by
the preceding estimates, and
$x\partial_x\Li_s(-zx)=\Li_{s-1}(-zx)$. Rewriting
$x/(1+x)=1-(1+x)^{-1}$ gives the formula.
\end{proof}
This recurrence is useful, but a termwise continuation beyond its common
strip can obscure cancellations. The compact-interval recurrence in
Section~\ref{bmh:sec:kernels} has no such ambiguity at integer $a$.

\subsection{A finite alphabet and a weight refinement}
Use $t=x/(1+x)$, so
\begin{equation}\label{bmh:eq:compact-kernel}
 \frac{x^{a-1}\dd x}{(1+x)^N}
 =R_{N,a}(t)\dd t,
 \qquad R_{N,a}(t)=t^{a-1}(1-t)^{N-a-1}.
\end{equation}
Let $G(\alpha_1,\ldots,\alpha_k;t)$ be the iterated integral with kernels
$\dd t/(t-\alpha)$, with the usual logarithmic value for all-zero words.
For $z\ne1$ put $\lambda_z=(1-z)^{-1}$. Direct differentiation and the
value at zero give
\begin{equation}\label{bmh:eq:Landen-G}
 \Li_m\!\left(-\frac{zt}{1-t}\right)
 =G\big((e_0-e_1)^{m-1}(e_1-e_{\lambda_z});t\big).
\end{equation}
Words and powers on the right are noncommutative, and $G$ is extended
linearly. For $z=1$ omit $e_{\lambda_z}$. Indeed the order-one identity is
$\log(1-t)-\log(1-(1-z)t)$, and higher orders follow from
$\dd\log(t/(1-t))=\dd t/t-\dd t/(t-1)$.

For fixed positive scales, let $\bmhH_{\le d}$ be the vector space over
$\Q(z_1,\ldots,z_r)$ generated by finite endpoint values of iterated
integrals in the alphabet
\begin{equation}\label{bmh:eq:alphabet}
 \{0,1,\lambda_{z_1},\ldots,\lambda_{z_r}\}
\end{equation}
of word length at most $d$, including products whose total length is
at most $d$. Endpoint logarithms are combined before taking finite
limits. This is a \emph{representation filtration}; no assertion about
independence of numerical periods is built into its definition.

\begin{theorem}[Multiscale finite closure and residue-weight law]
\label{bmh:thm:closure}
Let $s_1,\ldots,s_r$ be positive integers, $z_j>0$, $N\ge1$, and let
$a$ be an integer with $1-r\le a\le N-1$. Put $V=\sum_j s_j$.
Every logarithmic moment of $\bmhM_N$ is a finite multiple-polylogarithm
expression in the alphabet~\eqref{bmh:eq:alphabet}, and its anchored
primitive has a constructive hyperlogarithmic expression. More precisely,
for $h\ge0$ the moment belongs to $\bmhH_{\le V+h+1}$, and
\begin{equation}\label{bmh:eq:top-weight}
 \partial_a^h\bmhM_N(a)-c_{N,a}\,
 \left.\partial_b^h\bmhM_1(b)\right|_{b=0}
 \ \in\ \bmhH_{\le V+h},
\end{equation}
where
\begin{equation}\label{bmh:eq:residue}
 c_{N,a}=\begin{cases}
 0,&a\ge1,\\
 (-1)^k\binom{N+k-1}{k},&a=-k\le0.
 \end{cases}
\end{equation}
The factor parameters on both sides of~\eqref{bmh:eq:top-weight} are the
same. Coincident scales, including $z_j=1$, are interpreted by their
convergent limits.
\end{theorem}
\begin{proof}
Since $a,N$ are integers in the stated range, $R_{N,a}$ is a Laurent
polynomial, with no pole at $t=1$. Products in~\eqref{bmh:eq:Landen-G}
are expanded by shuffle, and
$\log x=\log t-\log(1-t)$ is a weight-one hyperlogarithm.
For completeness, the integration algorithm is the following finite
induction. Decompose a rational prefactor into a derivative of a rational
function and simple-pole terms. A simple pole appends one integration
letter. Integrate the rational derivative by parts; differentiating an
iterated integral removes its first letter and lowers word length by
one. Repeating terminates at weight zero. Polynomial prefactors are
handled by the same integration-by-parts step. All poles involved are
in~\eqref{bmh:eq:alphabet}, so no new letter is introduced. This is the
standard rational-hyperlogarithm mechanism~\cite[Section 2.2]{Panzer},
with the alphabet here explicitly determined.

For the sharper assertion, write
\[
 R_{N,a}(t)=\frac{c_{N,a}}t+S'(t),\qquad S(1)=0.
\]
Extracting the coefficient of $t^{-1}$ gives~\eqref{bmh:eq:residue}.
Let $P(t)$ be the polylogarithm product multiplied by
$\log^h(t/(1-t))$. Its weight is at most $V+h$. The $c_{N,a}/t$
term is precisely the second term of~\eqref{bmh:eq:top-weight} and can
raise weight by one. For the remainder,
\[
 \int_0^1S'(t)P(t)\dd t=-\int_0^1S(t)P'(t)\dd t.
\]
There are no boundary terms: at one, $S(t)=O(1-t)$ while $P$ grows at
most logarithmically; at zero, if $a=-k$, then $S=O(t^{-k})$ and
$P=O(t^r(1+|\log t|)^h)$, with $k<r$. The case $a\ge0$ is easier.
The derivative $P'$ has weight at most $V+h-1$. Rational integration
can raise it by at most one, proving~\eqref{bmh:eq:top-weight}.

These manipulations may first be performed on $[\varepsilon,1-\varepsilon]$.
The original integrability and the displayed boundary estimates determine
the combined endpoint limits, including cancellations between individual
hyperlogarithms. In particular a separated divergent endpoint value is
never silently treated as an ordinary integral. The same algorithm with
a variable upper endpoint gives an antiderivative; its value at zero
fixes the anchoring constant. Dominated convergence in the original
integral gives the asserted coincident-scale limits.
\end{proof}

\begin{remark}[What this settles]
For two positive integer orders, the theorem answers the finite-closure
part of the repository question on its entire integer product strip,
including $a=-1$. It gives an upper bound on representation complexity,
not a lower bound on arithmetic depth. Nor does it turn derivatives in
a noninteger spectral order into ordinary finite-weight hyperlogarithms.
\end{remark}

\section{A connected double-Hurwitz master formula}\label{bmh:sec:master}

\subsection{Removing the leading-index-one divergences}
Define the pole-cancelled Hurwitz difference
\begin{equation}\label{bmh:eq:D}
 D_k(q)=\begin{cases}
 \zeta(k;q)-\zeta(k),&k\ge2,\\
 \psi(1)-\psi(q),&k=1.
 \end{cases}
\end{equation}
The second line is the value at $s=1$ of
$D_s(q)=\zeta(s;q)-\zeta(s)$; the two simple poles have the same residue.

\begin{definition}[Connected coordinates]\label{bmh:def:connected}
For positive integers $A,B$ and $\Rea q>0$, put
\begin{align}
 \bmhC_{A,B}(q)
 &=\zeta(A,B)-Z_q(A,B)+\zeta(A;q)D_B(q),
 && A\ge2,\ B\ge1,\label{bmh:eq:Cbulk}\\
 \bmhC_{1,B}(q)
 &=Z_q(B,1)-\zeta(B,1)+D_{B+1}(q)-D_1(q)\zeta(B),
 && B\ge2,\label{bmh:eq:Cedge}\\
 \bmhC_{1,1}(q)&=\frac12\bigl(D_1(q)^2+D_2(q)\bigr).
 \label{bmh:eq:Ccorner}
\end{align}
Only convergent double-zeta coordinates appear in these formulas.
\end{definition}

\begin{lemma}[Connected stuffle, shift, and harmonic values]
\label{bmh:lem:Cidentities}
The functions in Definition~\ref{bmh:def:connected} are holomorphic for
$\Rea q>0$ and satisfy
\begin{align}
 \bmhC_{A,B}(q)+\bmhC_{B,A}(q)
 &=D_A(q)D_B(q)+D_{A+B}(q),\label{bmh:eq:Cstuffle}\\
 \bmhC_{A,B}(q+1)-\bmhC_{A,B}(q)
 &=-q^{-A}D_B(q).\label{bmh:eq:Cshift}
\end{align}
In particular $\bmhC_{A,B}(1)=\bmhC_{A,B}(2)=0$. For every integer $L\ge1$,
\begin{equation}\label{bmh:eq:Charmonic}
 \bmhC_{A,B}(L)=\sum_{1\le k<j\le L-1}\frac1{j^A k^B}.
\end{equation}
\end{lemma}
\begin{proof}
Holomorphy follows from normal convergence and~\eqref{bmh:eq:D}.
For $A,B\ge2$, use the absolutely convergent stuffle identity
\[
 \zeta(A;q)\zeta(B;q)
 =Z_q(A,B)+Z_q(B,A)+\zeta(A+B;q)
\]
at $q$ and at one. This gives~\eqref{bmh:eq:Cstuffle}. At the edges it
follows directly from~\eqref{bmh:eq:Cedge}; at the corner it is the
definition.

For $A\ge2$ we have
\[
 Z_{q+1}(A,B)=Z_q(A,B)-q^{-B}\bigl(\zeta(A;q)-q^{-A}\bigr),
 \qquad D_B(q+1)=D_B(q)-q^{-B}.
\]
Substitute these and $\zeta(A,q+1)=\zeta(A;q)-q^{-A}$ into
\eqref{bmh:eq:Cbulk}. Every term cancels except $-q^{-A}D_B(q)$.
For $A=1,B\ge2$, apply the same shifts to~\eqref{bmh:eq:Cedge};
for $A=B=1$ expand the square in~\eqref{bmh:eq:Ccorner}.
This proves~\eqref{bmh:eq:Cshift} in every case. Since $D_B(1)=0$,
the values at one and two vanish. Finally
$D_B(j)=-H_{j-1}^{(B)}$, also for $B=1$. Iterating the shift gives
\eqref{bmh:eq:Charmonic}.
\end{proof}

\subsection{The finite transform}
For $W=m+n$ define
\begin{align}
 \bmhK_{m,n}(q)={}&
 \sum_{k=0}^{m-1}\binom{W-2-k}{n-1}
       \bmhC_{W-1-k,k+1}(q)\notag\\
 &+\sum_{k=0}^{n-1}\binom{W-2-k}{m-1}
       \bmhC_{W-1-k,k+1}(q).
 \label{bmh:eq:K}
\end{align}
There are $m+n$ displayed terms, with coincident coordinates combinable.

\begin{theorem}[Bilinear Mellin--Hurwitz transform]
\label{bmh:thm:master}
For positive integers $m,n$ and $-2<\Rea a<1$, the ordinary integral
$M_1(a;m,n)$ satisfies
\begin{equation}\label{bmh:eq:master}
 \boxed{\displaystyle
 M_1(a;m,n)=\frac{\pi}{\sin\pi a}\,\bmhK_{m,n}(1-a).}
\end{equation}
At $a=0,-1$ the right side is interpreted by its removable value.
Every fixed derivative with respect to $a$ gives the corresponding
ordinary logarithmic moment. The formula requires no divergent value
$\zeta(1,B)$ or $Z_q(1,B)$.
\end{theorem}

\subsection{Proof by a Fermi integral and two cones}
For positive integer $m$,
\begin{equation}\label{bmh:eq:Fermi}
 \Li_m(-x)=-\frac{x}{\Gamma(m)}\int_0^\infty
             \frac{u^{m-1}}{e^u+x}\dd u.
\end{equation}
For real $0<a<1$, the product of the two negative signs is positive,
and Tonelli's theorem gives a double integral whose inner kernel is
\[
 I_a(c,d)=\int_0^\infty
       \frac{x^{a+1}}{(1+x)(x+c)(x+d)}\dd x,
 \qquad c=e^u,\ d=e^v.
\]
For distinct $1,c,d$, partial fractions give
\begin{equation}\label{bmh:eq:resolvent}
 I_a(c,d)=\frac\pi{\sin\pi a}
 \sum_{\rho\in\{1,c,d\}}
       \frac{\rho^{a+1}}{\prod_{\sigma\ne\rho}(\rho-\sigma)}.
\end{equation}
One first integrates the separate partial fractions in
$-2<\Rea a<-1$, where each is integrable, and then continues to the
full strip $-2<\Rea a<1$. The original three-denominator integral is
holomorphic there. Coincident $c,d,1$ are handled by continuity; this is
a divided difference, not a genuine pole at a collision.

Split the $(u,v)$ quadrant into $u>v$ and $v>u$. On the first cone set
$u=v+t$, $\xi=e^{-v}$, $\eta=e^{-t}$, and $q=1-a$. The bracket in
\eqref{bmh:eq:resolvent} becomes
\begin{equation}\label{bmh:eq:cone-kernel}
 \mathscr B_a(v,t)=
 \frac{\xi^2\eta}{(1-\xi\eta)(1-\xi)}
 +\frac{(\xi\eta)^q}{(1-\xi\eta)(1-\eta)}
 -\frac{\xi^q\eta}{(1-\xi)(1-\eta)}.
\end{equation}
The identity follows by substituting $c=(\xi\eta)^{-1}$ and
$d=\xi^{-1}$ into the three rational coefficients in
\eqref{bmh:eq:resolvent}.

We require its two-variable Mellin moments. For $\Rea A,\Rea B>1$,
expand the three geometric kernels and integrate separately. The result is
\begin{align}
 \frac1{\Gamma(A)\Gamma(B)}\int_0^\infty\!\int_0^\infty
 v^{A-1}t^{B-1}\mathscr B_a(v,t)\dd v\dd t
 &=\zeta(A,B)+Z_q(B,A)+\zeta(A+B;q)\notag\\
 &\quad-\zeta(A;q)\zeta(B)
 =\bmhC_{A,B}(q).
 \label{bmh:eq:cone-moment}
\end{align}
For example, the first geometric expansion has frequencies
$(j+k+2,j+1)$, with the first strictly greater than the second, and
therefore produces $\zeta(A,B)$. The second has frequencies
$(q+j,q+j+k)$ and produces a weakly ordered double sum. Its diagonal is
$\zeta(A+B;q)$. The third factorizes. The last equality is precisely the
stuffle cancellation.

It remains to justify the boundary cases $A=1$ or $B=1$, rather than
subtract divergent series informally. For fixed real $0<a<1$,
$\mathscr B_a$ is bounded near the origin: it is
$(\sin\pi a/\pi)I_a(e^{v+t},e^v)$, and the integral is smooth when its
three poles coalesce off the integration path. At infinity it decays
exponentially in $v+t$. Indeed, replacing $1+x$ by $x$ gives the upper
bound
\[
 I_a(e^{v+t},e^v)\le
 \frac\pi{\sin\pi a}\,e^{-(1-a)v}
       \frac{e^{at}-1}{e^t-1}.
\]
Thus the combined integral in~\eqref{bmh:eq:cone-moment} is holomorphic
for $\Rea A,\Rea B>0$.

For $A>1$, the limit $B\to1$ on the right is
\eqref{bmh:eq:Cbulk}, because the common residue has cancelled in
$D_B$. Using~\eqref{bmh:eq:Cstuffle} gives the limit $A\to1,B>1$ as
\eqref{bmh:eq:Cedge}. On the diagonal $A=B\to1$,
\eqref{bmh:eq:Cstuffle} gives~\eqref{bmh:eq:Ccorner}. Holomorphy of the
combined integral proves that these are its genuine boundary values.
This argument explicitly distinguishes cancellation in the integrand
from nonexistent separate integrals at the leading-one coordinates.

Finally expand
\[
 (v+t)^{m-1}v^{n-1}
 =\sum_{k=0}^{m-1}\binom{m-1}{k}
       v^{W-2-k}t^k.
\]
After the Gamma normalizations in~\eqref{bmh:eq:Fermi} and
\eqref{bmh:eq:cone-moment}, the coefficient is
\[
 \frac{\binom{m-1}{k}\Gamma(W-1-k)\Gamma(k+1)}
      {\Gamma(m)\Gamma(n)}
 =\binom{W-2-k}{n-1}.
\]
The other cone interchanges $m,n$. This proves~\eqref{bmh:eq:master}
for real $0<a<1$. Both sides are holomorphic, apart from the displayed
cosecant poles, on $-2<\Rea a<1$; the identity theorem extends the result.
The original integral proves removability at $0,-1$, also visible from
$\bmhC(1)=\bmhC(2)=0$. Proposition~\ref{bmh:prop:strip} justifies all fixed
logarithmic derivatives and completes the proof.

\begin{example}[The logarithm-squared case]\label{bmh:ex:11master}
For $m=n=1$ the master is especially transparent:
\begin{equation}\label{bmh:eq:11master}
 M_1(a;1,1)=\pi\csc(\pi a)
 \left((\psi(1)-\psi(1-a))^2+\zeta(2;1-a)-\zeta(2)\right).
\end{equation}
This also follows by differentiating the beta integral twice with respect
to its denominator exponent. It is a normalization check, not a claim of
priority for a classical beta derivative.
\end{example}

\begin{remark}[Derivatives away from resonance]
Every $q$-derivative of the double coordinates is finite:
\begin{equation}\label{bmh:eq:Zq-derivative}
 \frac{\dd^r}{\dd q^r}Z_q(A,B)=(-1)^r
 \sum_{j=0}^r\binom rj(A)_j(B)_{r-j}
       Z_q(A+j,B+r-j).
\end{equation}
Together with the derivative ladder for $D_k$ and the cosecant factor,
this gives explicit finite formulas for all logarithmic moments at
noninteger $a$. At a resonance one should instead use the cancelled
Taylor formulas in the next section.
\end{remark}

\section{Both resonances and every logarithmic moment}\label{bmh:sec:resonances}

\subsection{An all-weight double-zeta evaluation}
Write
\begin{equation}\label{bmh:eq:Idef}
 I_{m,n}^{(h)}=\int_0^\infty
 \frac{\log^h x\,\Li_m(-x)\Li_n(-x)}{x(1+x)}\dd x,
 \qquad I_{m,n}=I_{m,n}^{(0)}.
\end{equation}

\begin{theorem}[Compact double-zeta reduction]\label{bmh:thm:compact}
For every $m,n\ge1$, with $W=m+n$,
\begin{align}
 \boxed{\begin{aligned}
 I_{m,n}={}&\left(\binom W{m-1}+\binom W{n-1}\right)\zeta(W+1)\\
 &+\sum_{j=1}^{W-1}j\,A_j(m,n)\zeta(j+1,W-j).
 \end{aligned}}
 \label{bmh:eq:compact}
\end{align}
where
\begin{equation}\label{bmh:eq:Aj}
 A_j(m,n)=\binom{W-j-1}{m-1}+\binom{W-j-1}{n-1}
           -\binom{j-1}{m-1}-\binom{j-1}{n-1}.
\end{equation}
In particular $A_{W-j}(m,n)=-A_j(m,n)$. All double-zeta values in the
formula are ordinarily convergent and have weight $W+1$.
\end{theorem}

\begin{proof}[A direct proof at resonance]
The $a=0$ limit of~\eqref{bmh:eq:resolvent} is
\[
 I_0(c,d)=\frac{c\log c/(c-1)-d\log d/(d-1)}{c-d}.
\]
On the cone $u=v+t$ it is the following difference of two nonnegative
kernels:
\begin{equation}\label{bmh:eq:resonant-cone}
 I_0(e^{v+t},e^v)=
 \frac{t e^{-(v+t)}}{(1-e^{-t})(1-e^{-(v+t)})}
 -\frac{v e^{-(2v+t)}}{(1-e^{-(v+t)})(1-e^{-v})}.
\end{equation}
Each term, after multiplication by $(v+t)^{m-1}v^{n-1}$, is integrable.
This can be seen either from its elementary behavior at the axes or from
the convergent double sums below. Consequently Tonelli's theorem applies
separately to their geometric expansions.

Using the same binomial expansion as in the proof of the master, the
first cone contributes
\begin{align}
 \sum_{k=0}^{m-1}\binom{W-2-k}{n-1}
 \bigl[&(k+1)\{\zeta(k+2,W-1-k)+\zeta(W+1)\}\notag\\
       &-(W-1-k)\zeta(W-k,k+1)\bigr].
 \label{bmh:eq:resonant-unsimplified}
\end{align}
The weak inequality in the first geometric sum accounts for the
single-zeta diagonal. The second cone gives the same expression with
$m,n$ interchanged. Reindexing the positive and negative double-zeta
terms gives $jA_j(m,n)$. The coefficient of the diagonal is
\[
 \sum_{k=0}^{m-1}(k+1)\binom{W-2-k}{n-1}
 +\sum_{k=0}^{n-1}(k+1)\binom{W-2-k}{m-1}
 =\binom W{m-1}+\binom W{n-1},
\]
a double application of the hockey-stick identity. This proves the
formula independently of any numerical recognition.
\end{proof}

\subsection{All Taylor coefficients without singular coordinates}
Let
\begin{equation}\label{bmh:eq:even-coeff}
 \frac{\pi u}{\sin\pi u}=\sum_{j\ge0}e_{2j}u^{2j},
 \qquad e_0=1,\quad
 e_{2j}=2(1-2^{1-2j})\zeta(2j)\quad(j\ge1).
\end{equation}
For positive integers $A,B$ and $\nu\ge1$, define
\begin{align}
 c_\nu(A,B)={}&-\frac{(A)_\nu}{\nu!}\zeta(A+\nu,B)\notag\\
 &+\sum_{v=1}^{\nu}
 \frac{(A)_{\nu-v}(B)_v}{(\nu-v)!v!}
 \bigl\{\zeta(B+v,A+\nu-v)+\zeta(A+B+\nu)\bigr\}.
 \label{bmh:eq:cjet}
\end{align}
Every first double-zeta index in this expression is at least two. An
equivalent, simpler corner formula is
\begin{equation}\label{bmh:eq:c11jet}
 c_\nu(1,1)=\frac12\sum_{v=1}^{\nu-1}
       \zeta(v+1)\zeta(\nu-v+1)
       +\frac{\nu+1}{2}\zeta(\nu+2).
\end{equation}
Define $T_\nu(m,n)$ by replacing every $\bmhC_{A,B}$ in
\eqref{bmh:eq:K} by $c_\nu(A,B)$.

\begin{lemma}[Cancelled coefficient identity]\label{bmh:lem:cjet}
For $|u|<1$,
\begin{equation}\label{bmh:eq:C-Taylor}
 \bmhC_{A,B}(1-u)=\sum_{\nu\ge1}c_\nu(A,B)u^\nu.
\end{equation}
\end{lemma}
\begin{proof}
For $A\ge2$, expand the normally convergent double Hurwitz series and
use
\[
 D_B(1-u)=\sum_{v\ge1}\frac{(B)_v}{v!}\zeta(B+v)u^v.
\]
The latter also holds at $B=1$ by the digamma expansion. At total degree
$\nu$, the terms of $Z_{1-u}(A,B)$ with positive shift in the second
index cancel against the corresponding stuffle term of
$\zeta(A;1-u)D_B(1-u)$. What remains is exactly
\eqref{bmh:eq:cjet}. For $A=1$, take the limit from spectral $A>1$ in
the analytic combined cone integral of~\eqref{bmh:eq:cone-moment}.
All coordinates surviving on the right of~\eqref{bmh:eq:cjet} converge
at this limit. The same argument reaches $A=B=1$; alternatively,
expanding~\eqref{bmh:eq:Ccorner} directly gives
\eqref{bmh:eq:c11jet}. Local uniform convergence permits coefficient
extraction throughout.
\end{proof}

\begin{theorem}[Complete resonant logarithmic moments]
\label{bmh:thm:jets}
For every $m,n\ge1$ and $h\ge0$,
\begin{equation}\label{bmh:eq:all-jets}
 \boxed{\displaystyle
 I_{m,n}^{(h)}=h!\sum_{j=0}^{\lfloor h/2\rfloor}
                      e_{2j}T_{h+1-2j}(m,n).}
\end{equation}
Every term has representation weight $m+n+h+1$. At $h=0$ the formula
reduces to Theorem~\ref{bmh:thm:compact}.
\end{theorem}
\begin{proof}
Insert~\eqref{bmh:eq:C-Taylor} into~\eqref{bmh:eq:master} and use
\eqref{bmh:eq:even-coeff}. The constant term of $\bmhK(1-u)$ is zero, so
its factor of $u$ cancels the cosecant pole. The coefficient of $u^h$
is the right side of~\eqref{bmh:eq:all-jets} divided by $h!$.
Differentiating the original integral is justified by
Proposition~\ref{bmh:prop:strip}. The weight assertion follows by
adding the indices in~\eqref{bmh:eq:cjet} and the even-zeta weight.
\end{proof}

\begin{remark}[Depth is not the same as the generating algebra]
For $h=0,1$ the answer is a linear combination of single and double zeta
values. For larger $h$, the factor $e_{2j}$ may multiply a double value.
It is therefore correct to say that the answer belongs to the algebra
generated by ordinary and double zeta values. It is \emph{not} correct
to call every such product additive depth two. Expanding by stuffle gives
additive depth at most three. No minimal-depth assertion is made.
\end{remark}

\subsection{The second resonance at \texorpdfstring{$a=-1$}{a=-1}}
Define
\begin{equation}\label{bmh:eq:minus-coeff}
 c^-_\nu(A,B)=c_\nu(A,B)
 -\sum_{v=1}^{\nu}
    \frac{(A)_{\nu-v}(B)_v}{(\nu-v)!v!}\zeta(B+v),
\end{equation}
and let $T^-_\nu(m,n)$ be the corresponding replacement in
\eqref{bmh:eq:K}.

\begin{corollary}[Negative Mellin resonance]\label{bmh:cor:negative}
For $m,n\ge1$ and $h\ge0$,
\begin{equation}\label{bmh:eq:negative-jets}
 \boxed{\displaystyle
 \int_0^\infty\frac{\log^h x\,\Li_m(-x)\Li_n(-x)}{x^2(1+x)}\dd x
 =-h!\sum_{j=0}^{\lfloor h/2\rfloor}e_{2j}T^-_{h+1-2j}(m,n).}
\end{equation}
\end{corollary}
\begin{proof}
The shift law gives
\[
 \bmhC_{A,B}(2-u)=\bmhC_{A,B}(1-u)-(1-u)^{-A}D_B(1-u).
\]
Multiplying the two elementary Taylor series proves that its coefficients
are $c^-_\nu$. Now use $a=-1+u$ and
$\sin\pi(-1+u)=-\sin\pi u$ in~\eqref{bmh:eq:master}.
\end{proof}
For instance,
\begin{equation}\label{bmh:eq:negative11}
 \int_0^\infty\frac{\log^2(1+x)}{x^2(1+x)}\dd x
 =2\zeta(2)-2\zeta(3).
\end{equation}
This resonance exists because the product is $O(x^2)$ at zero. The
one-factor lower endpoint $a>-1$ cannot be reused for the bilinear
problem.

\subsection{Small-weight reductions with explicit algebra}
The first values are
\begin{align}
 I_{1,1}&=2\zeta(3),\label{bmh:eq:example11}\\
 I_{2,1}&=\frac{17}{4}\zeta(4),\label{bmh:eq:example21}\\
 I_{3,1}&=\frac{11}{2}\zeta(5)+\zeta(2)\zeta(3),\label{bmh:eq:example31}\\
 I_{2,2}&=2\zeta(5)+4\zeta(2)\zeta(3).\label{bmh:eq:example22}
\end{align}
Here is enough exact algebra to verify the reductions. The depth-two sum
formula
\begin{equation}\label{bmh:eq:sumformula}
 \sum_{k=2}^{w-1}\zeta(k,w-k)=\zeta(w),\qquad w\ge3,
\end{equation}
follows without regularization: sum the finite geometric progression
inside $\sum_{n>m}$ to obtain
\[
 \sum_{k=2}^{w-1}\frac1{n^k m^{w-k}}
 =\frac{1/(n m^{w-2})-1/n^{w-1}}{n-m}.
\]
The first part summed over $n>m$ is
$\sum_m H_m/m^{w-1}$, and the second is
$\sum_n H_{n-1}/n^{w-1}$. Their difference is $\zeta(w)$.
At weight four, stuffle gives $\zeta(2,2)=3\zeta(4)/4$, so
$\zeta(3,1)=\zeta(4)/4$. At weight five, use the convergent shuffle
and stuffle rows
\begin{align*}
 \zeta(2)\zeta(3)&=\zeta(2,3)+3\zeta(3,2)+6\zeta(4,1),\\
 \zeta(2)\zeta(3)&=\zeta(2,3)+\zeta(3,2)+\zeta(5),
\end{align*}
together with~\eqref{bmh:eq:sumformula}. Solving gives
\[
 \zeta(4,1)=2\zeta(5)-\zeta(2)\zeta(3),\quad
 \zeta(2,3)=\frac92\zeta(5)-2\zeta(2)\zeta(3).
\]
Substitution into~\eqref{bmh:eq:compact} proves the displayed examples.

\section{All integer kernels: a terminating harmonic recurrence}
\label{bmh:sec:kernels}

\subsection{Laurent transport on the compact interval}
Put
\[
 y(t)=\log\frac{t}{1-t},\qquad
 L_j(t)=\Li_j\!\left(-\frac{t}{1-t}\right).
\]
The identities required by the recurrence are
\begin{equation}\label{bmh:eq:Lderivative}
 L_0(t)=-t,\qquad
 L_j'(t)=\frac{L_{j-1}(t)}{t(1-t)}\quad(j\ge1),\qquad
 y'(t)=\frac1{t(1-t)}.
\end{equation}
In particular the zero-order value is $-t$, not $-t/(1-t)$.

For a Laurent polynomial $R(t)=\sum_jr_jt^j$, define
\begin{equation}\label{bmh:eq:SR}
 c_R=r_{-1},\qquad
 S_R(t)=\sum_{j\ne-1}r_j\frac{t^{j+1}-1}{j+1},\qquad
 K_R(t)=\frac{S_R(t)}{t(1-t)}.
\end{equation}
Since $S_R(1)=0$, $K_R$ is again a Laurent polynomial. Let
\[
 F_R(m,n;h)=\int_0^1 R(t)y(t)^h L_m(t)L_n(t)\dd t,
\]
whenever the displayed integral is ordinary and convergent.

\begin{theorem}[Finite integer-kernel reduction]\label{bmh:thm:kernel}
For $m,n\ge1$ and every integrable instance reached from
$R=R_{N,a}$ with $-1\le a\le N-1$ integer,
\begin{align}
 F_R(m,n;h)={}&c_R I_{m,n}^{(h)}
 -F_{K_R}(m-1,n;h)-F_{K_R}(m,n-1;h)\notag\\
 &-hF_{K_R}(m,n;h-1),\label{bmh:eq:kernel-rec}
\end{align}
where the last term is omitted when $h=0$. Whenever an index reaches
zero, absorb $L_0=-t$ into the rational kernel and remove that factor.
The one-factor recurrence is
\begin{equation}\label{bmh:eq:one-rec}
 F_R(n;h)=c_R J_n^{(h)}-F_{K_R}(n-1;h)-hF_{K_R}(n;h-1),
\end{equation}
with
\begin{equation}\label{bmh:eq:one-base}
 J_n^{(h)}=-h!\sum_{j=0}^{\lfloor h/2\rfloor}
 e_{2j}\frac{(n)_{h+1-2j}}{(h+1-2j)!}\,
            \zeta(n+h+1-2j).
\end{equation}
The final zero-factor bases are the Gamma--harmonic moments in
Theorem~\ref{bmh:thm:beta-base}. These rules terminate and evaluate every
$\partial_a^h M_N(a;m,n)$ at every admissible integer $a$.
\end{theorem}
\begin{proof}
Write $R=c_R/t+S_R'$ and integrate the second term by parts. Use
\eqref{bmh:eq:Lderivative} on each factor and on $y^h$ to get
\eqref{bmh:eq:kernel-rec}. The boundary at one vanishes because
$S_R(1)=0$ and each $L_j$ and $y$ is at most logarithmic there. At zero,
the initial bilinear kernel has pole order at most two. Two positive-order
factors are $O(t^2)$, and $S_R$ has pole order at most one, so the boundary
vanishes even with logarithmic powers. The recurrence preserves this
property. Absorbing a zero-order factor multiplies the kernel by $-t$;
with one factor left the pole order is at most one, and with no factors
left the kernel is an ordinary polynomial. Thus all reached boundary
terms and integrals have the asserted ordinary meaning.

The same argument proves~\eqref{bmh:eq:one-rec}. Its base
$J_n^{(h)}=\int_0^1y^hL_n(t)\dd t/t$ is the one-factor resonance from
\cite{MLD}. To check its normalization directly, use the already proved
one-factor transform
\[
 J_1(a,n;1)=\pi\csc(\pi a)\{\zeta(n)-\zeta(n;1-a)\}
\]
and its cancelled Taylor series at $a=0$. At $n=1$, replace the zeta
difference by $\psi(1-a)-\psi(1)$ before expansion. This gives exactly
\eqref{bmh:eq:one-base}, including its minus sign.

Every recursive call reduces the sum of the positive orders and $h$ by
one, or removes a zero-order factor. Polynomial decomposition is finite.
The process therefore terminates without solving an infinite recurrence.
Finally~\eqref{bmh:eq:compact-kernel} identifies the initial integral
with the desired Mellin moment.
\end{proof}

\begin{corollary}[Representation class at integer parameters]
At $h=0$, every unit-scale $M_N(a;m,n)$ with $a=-1,0,\ldots,N-1$ is a
rational linear combination of rational numbers, ordinary zeta values,
and convergent double zeta values, of weights at most $m+n+1$.
For $a\ge1$ the weight bound improves to $m+n$. With $h$ logarithms the
bounds become $m+n+h+1$ and $m+n+h$, respectively. The answers belong
to the algebra generated by single and double zeta values; after stuffle,
an additive depth bound of three is sufficient for the formulas produced
here.
\end{corollary}
\begin{proof}
Use Theorem~\ref{bmh:thm:jets}, the one-factor base, and the elementary
bases below in the terminating recurrence. The sharper weight bounds
are also immediate from~\eqref{bmh:eq:top-weight}. In the coefficient
formulas the only extra products are even zeta factors multiplying a
single or double value, or the two-zeta corner terms. Products of even
zeta values are rational multiples of a single even zeta value. Thus
stuffle requires at most three nested indices.
\end{proof}

\subsection{The Gamma base and generalized harmonic numbers}
\begin{theorem}[Beta moments and Bell polynomials]\label{bmh:thm:beta-base}
For integers $d,h\ge0$,
\begin{align}
 B_{d,h}&:=\int_0^1t^d\log^h\frac{t}{1-t}\dd t\notag\\
 &=\left.\frac{\dd^h}{\dd v^h}B(d+1+v,1-v)\right|_{v=0}\notag\\
 &=\frac{h!}{d+1}[v^h]\left\{
       \frac{\pi v}{\sin\pi v}\prod_{j=1}^d\left(1+\frac vj\right)
       \right\}.\label{bmh:eq:beta-base}
\end{align}
Equivalently, $B_{d,h}=Y_h(\kappa_1,\ldots,\kappa_h)/(d+1)$, where
$Y_h$ is the complete exponential Bell polynomial and
\begin{align}
 \kappa_1&=H_d,\label{bmh:eq:kappa1}\\
 \kappa_r&=(r-1)!\left\{
 (1+(-1)^r)\zeta(r)+(-1)^{r+1}H_d^{(r)}\right\},\qquad r\ge2.
 \label{bmh:eq:kappas}
\end{align}
\end{theorem}
\begin{proof}
Differentiate the ordinary beta integral near $v=0$; an integrable
power-logarithmic majorant justifies the derivatives. The Gamma quotient
and reflection identity give
\[
 B(d+1+v,1-v)
 =\frac1{d+1}\Gamma(1+v)\Gamma(1-v)
           \prod_{j=1}^d(1+v/j)
 =\frac1{d+1}\frac{\pi v}{\sin\pi v}\prod_{j=1}^d(1+v/j).
\]
This proves~\eqref{bmh:eq:beta-base}. Taking logarithmic derivatives
of the product gives~\eqref{bmh:eq:kappa1}--\eqref{bmh:eq:kappas}.
The relation between derivatives of an exponential and the complete
Bell polynomials proves the last assertion. These beta and Gamma tools
are classical; their role here is to close the bilinear recurrence.
\end{proof}

\subsection{Examples with different denominator orders}
The recurrence gives, without numerical fitting,
\begin{align}
 M_2(1;1,1)&=2,\label{bmh:eq:kernel-example1}\\
 M_3(1;2,1)&=-\frac54+\frac14\zeta(2)+\zeta(3),\label{bmh:eq:kernel-example2}\\
 M_3(2;2,2)&=\zeta(2)+2\zeta(3)+\frac{17}{4}\zeta(4),\label{bmh:eq:kernel-example3}\\
 M_2(0;2,1)&=\frac{17}{4}\zeta(4)-\zeta(2)-2\zeta(3),\label{bmh:eq:kernel-example4}\\
 \left.\partial_a M_2(a;2,2)\right|_{a=1}
 &=15\zeta(5)+12\zeta(2)\zeta(3).\label{bmh:eq:kernel-example5}
\end{align}
Their exact expressions are regenerated by the supplied symbolic code.
The two-factor weight bound does not imply homogeneity when the rational
kernel has polynomial parts: the lower-weight terms in these examples
are expected, not anomalous.

\section{Explicit primitives and a weighted multiple-zeta identity}
\label{bmh:sec:words}

\subsection{Binary Landen words}
In this section it is convenient to use the positive harmonic kernels
$\omega_0=\dd t/t$ and $\omega_1=\dd t/(1-t)$. Let $H(w;t)$ denote the
corresponding iterated integral, with the first letter integrated
outermost. These conventions differ by a sign at the letter one from
those of $G$ in Section~\ref{bmh:sec:closure}. Equation~\eqref{bmh:eq:Landen-G}
at $z=1$ becomes
\begin{equation}\label{bmh:eq:Landen-binary}
 L_m(t)=-H\big((e_0+e_1)^{m-1}e_1;t\big).
\end{equation}
For every word $w=w_1\cdots w_W$ of length $W=m+n$ ending in one, put
\begin{equation}\label{bmh:eq:word-coeff}
 c_{m,n}(w)=\sum_{\substack{1\le j<W\\w_j=1}}
 \left\{\binom{j-1}{m-1}+\binom{j-1}{n-1}\right\}.
\end{equation}

\begin{lemma}[Closed shuffle multiplicities]\label{bmh:lem:shuffle}
For $0<t<1$,
\begin{equation}\label{bmh:eq:shuffle-product}
 L_m(t)L_n(t)=\sum_{\substack{|w|=W\\w_W=1}}c_{m,n}(w)H(w;t).
\end{equation}
\end{lemma}
\begin{proof}
The two minus signs in~\eqref{bmh:eq:Landen-binary} cancel. A shuffle
assigns $m$ positions of $w$ to the first factor and $n$ to the second.
Every internal letter in either factor may be zero or one, but its final
letter must be one. Exactly one factor ends at position $W$ and the
other ends at some $j<W$ with $w_j=1$. If the length-$m$ factor ends at
$j$, choose its other $m-1$ positions among the first $j-1$ positions,
giving $\binom{j-1}{m-1}$ possibilities. Interchanging the factors gives
the second binomial coefficient. These cases are disjoint and exhaustive,
proving~\eqref{bmh:eq:word-coeff}.
\end{proof}

\begin{theorem}[An explicit anchored primitive]\label{bmh:thm:primitive}
For $x\ge0$,
\begin{equation}\label{bmh:eq:primitive}
 \boxed{\displaystyle
 \int_0^x\frac{\Li_m(-u)\Li_n(-u)}{u(1+u)}\dd u
 =\sum_{\substack{|w|=W\\w_W=1}}c_{m,n}(w)
       H\left(0w;\frac{x}{1+x}\right).}
\end{equation}
Every endpoint value at $x=\infty$ is convergent: its word starts with
zero and ends with one.
\end{theorem}
\begin{proof}
The substitution $t=u/(1+u)$ turns $\dd u/(u(1+u))$ into $\dd t/t$.
Insert~\eqref{bmh:eq:shuffle-product} and prepend the letter zero.
Both sides vanish at $x=0$, fixing the integration constant. At one the
prepend-zero condition removes the outer logarithmic singularity, while
the final-one condition gives convergence at zero.
\end{proof}

\subsection{A finite high-depth sum with a double-zeta answer}
Let $\boldsymbol k=(k_1,\ldots,k_d)$ be a composition of $W+1$ with
$k_1\ge2$, and write $K_\ell=k_1+\cdots+k_\ell$. Define
\begin{equation}\label{bmh:eq:composition-coeff}
 c_{m,n}(\boldsymbol k)=\sum_{\ell=1}^{d-1}
 \left\{\binom{K_\ell-2}{m-1}+\binom{K_\ell-2}{n-1}\right\}.
\end{equation}
The sum is zero when $d=1$.

\begin{corollary}[Weighted MZV reduction]\label{bmh:cor:weighted}
For all positive $m,n$,
\begin{align}
 &\sum_{\substack{k_1+\cdots+k_d=W+1\\k_1\ge2,\ k_j\ge1}}
 c_{m,n}(\boldsymbol k)\zeta(k_1,\ldots,k_d)\notag\\
 &\qquad=\left(\binom W{m-1}+\binom W{n-1}\right)\zeta(W+1)
 +\sum_{j=1}^{W-1}jA_j(m,n)\zeta(j+1,W-j).
 \label{bmh:eq:weighted}
\end{align}
The left side ranges over every admissible composition, and can have
large depth. Its particular weighted combination always reduces to
single and double values.
\end{corollary}
\begin{proof}
Take $x\to\infty$ in~\eqref{bmh:eq:primitive} and identify
$H(0^{k_1-1}1\cdots0^{k_d-1}1;1)$ with the decreasing-index multiple
zeta value. An internal one occurs at position $K_\ell-1$ in $w$, so
\eqref{bmh:eq:word-coeff} becomes~\eqref{bmh:eq:composition-coeff}.
Theorem~\ref{bmh:thm:compact} supplies the right side.
\end{proof}
This is a reduction of a specified family of weighted sums, not a theorem
that every multiple zeta value of the same weight has depth at most two.
It also gives a direct identity between two finite descriptions of the
same integral: a binary-word primitive and a two-cone double-sum formula.
The verification code checks their combinatorial interface without
assuming numerical period independence.

\section{Inversion reciprocity and central parity reductions}
\label{bmh:sec:reciprocity}

\subsection{The exact inversion polynomial}
For $n\ge1$ put
\begin{equation}\label{bmh:eq:Pn}
 P_n(y)=-\frac{y^n}{n!}
 -2\sum_{k=1}^{\lfloor n/2\rfloor}
 (1-2^{1-2k})\zeta(2k)\frac{y^{n-2k}}{(n-2k)!}.
\end{equation}
Thus $P_1=-y$, $P_2=-y^2/2-\zeta(2)$, and
$P_3=-y^3/6-\zeta(2)y$.

\begin{lemma}[Integer-order inversion]\label{bmh:lem:inversion}
For real $y$,
\begin{equation}\label{bmh:eq:inversion}
 \Li_n(-e^y)+(-1)^n\Li_n(-e^{-y})=P_n(y).
\end{equation}
Also $P_n'=P_{n-1}$ for $n\ge2$ and
$P_n(-y)=(-1)^nP_n(y)$.
\end{lemma}
\begin{proof}
The order-one statement is
$-\log(1+e^y)+\log(1+e^{-y})=-y$.
Differentiating the left side at order $n$ gives the left side at
order $n-1$. The same is true of the displayed polynomial.
At $y=0$, the value is zero for odd $n$, and is
$2\Li_n(-1)=-2(1-2^{1-n})\zeta(n)$ for even $n$.
These constants determine the induction. The last parity assertion is
visible term by term. This is the classical integer inversion identity;
the proof fixes the branch and normalization needed below.
\end{proof}

Let
\begin{equation}\label{bmh:eq:JandBeta}
 J_N(a,s)=\int_0^\infty\frac{x^{a-1}\Li_s(-x)}{(1+x)^N}\dd x,
 \qquad B_N(a)=B(a,N-a).
\end{equation}
If $P$ is a polynomial, $P(\partial_a)$ denotes its constant-coefficient
differential operator.

\begin{theorem}[Bilinear reciprocity]\label{bmh:thm:reciprocity}
For $0<\Rea a<N$,
\begin{align}
 M_N(a;m,n)-(-1)^W M_N(N-a;m,n)
 ={}&P_m(\partial_a)J_N(a,n)
    +P_n(\partial_a)J_N(a,m)\notag\\
 &-(P_mP_n)(\partial_a)B_N(a).
 \label{bmh:eq:reciprocity}
\end{align}
Every fixed logarithmic derivative is valid in the same strip.
\end{theorem}
\begin{proof}
Set $L_j(x)=\Li_j(-x)$. Inversion gives
$L_j(1/x)=(-1)^j(P_j(\log x)-L_j(x))$.
Multiply the two formulas to obtain
\[
 L_m(x)L_n(x)-(-1)^W L_m(1/x)L_n(1/x)
 =P_m(\log x)L_n(x)+P_n(\log x)L_m(x)-P_m(\log x)P_n(\log x).
\]
Integrate with the Mellin kernel. In the second bilinear term, $x\mapsto
1/x$ changes $a$ into $N-a$. All separated terms converge for
$0<\Rea a<N$; in particular the purely polynomial term is an ordinary
beta logarithmic moment. This proves~\eqref{bmh:eq:reciprocity}.
\end{proof}

\begin{theorem}[Complete odd-parity central moments]
\label{bmh:thm:central}
If $W+h$ is odd, then
\begin{align}
 &\int_0^\infty\frac{x^{N/2-1}\log^h x\,
                    \Li_m(-x)\Li_n(-x)}{(1+x)^N}\dd x\notag\\
 &\quad=\frac12\left.\partial_a^h
 \left\{P_m(\partial_a)J_N(a,n)+P_n(\partial_a)J_N(a,m)
 -(P_mP_n)(\partial_a)B_N(a)\right\}\right|_{a=N/2}.
 \label{bmh:eq:central}
\end{align}
For every integer $N$, the right side is a finite expression in rational
numbers, $\pi$, $\log2$, and ordinary positive-integer zeta values.
When $W+h$ is even, reciprocity yields a relation among the one-factor
terms but does not determine the central bilinear moment.
\end{theorem}
\begin{proof}
The $h$th derivative of the left side of~\eqref{bmh:eq:reciprocity} at
$a=N/2$ is $(1-(-1)^{W+h})$ times the desired integral. This is two
exactly in the stated parity sector.

To establish the expression class, use the existing finite one-factor
formula~\cite{MLD}
\begin{equation}\label{bmh:eq:priorJN}
 J_N(a,s)=(-1)^N\pi\csc(\pi a)
       \sum_{j=0}^{N-1}q_{N,j}(a)
              \{\zeta(s-j;N-a)-\zeta(s-j)\},
\end{equation}
where
\[
 \sum_jq_{N,j}(a)X^j=
   \frac1{(N-1)!}\prod_{r=1}^{N-1}(X+a-r).
\]
It is interpreted as a combined holomorphic germ at integer $a$ and at
spectral zeta poles. If $N$ is odd, $N/2$ is a half-integer. The shift
identity and
$\zeta(k;1/2)=(2^k-1)\zeta(k)$ reduce all regular Hurwitz values;
at $k=1$ the cancelled difference is a digamma difference and contributes
$\log2$ and rational numbers. Nonpositive integer values are Bernoulli
polynomials. The $a$ derivatives shift Hurwitz orders upward, with
removable products retained before evaluation. If $N$ is even, the
integer one-factor resonance and its logarithmic moments reduce to
\eqref{bmh:eq:one-base} through the finite kernel recurrence. Finally
$B_N$ is a Gamma quotient with integer or half-integer parameters, whose
logarithmic derivatives are the corresponding polygamma values. These
reduce by the same shift and half-argument formulas. This proves the
finite expression claim without a divergent termwise substitution.
\end{proof}

\subsection{Explicit central evaluations}
At $N=1$, the first two examples are
\begin{align}
 \int_0^\infty\frac{\Li_1(-x)\Li_2(-x)}{\sqrt{x}(1+x)}\dd x
 &=\frac{21}{2}\pi\zeta(3)+\frac23\pi^3\log2,
 \label{bmh:eq:central12}\\
 \int_0^\infty\frac{\log x\,\log^2(1+x)}{\sqrt{x}(1+x)}\dd x
 &=14\pi\zeta(3)+2\pi^3\log2.
 \label{bmh:eq:central11log}
\end{align}
Further exact instances are
\begin{align}
 \int_0^\infty\frac{\Li_2(-x)\Li_3(-x)}{\sqrt{x}(1+x)}\dd x
 &=155\pi\zeta(5)+\frac{20}{3}\pi^3\zeta(3),
 \label{bmh:eq:central23}\\
 \int_0^\infty\frac{\log x\,\Li_2(-x)^2}{\sqrt{x}(1+x)}\dd x
 &=372\pi\zeta(5)+\frac{70}{3}\pi^3\zeta(3).
 \label{bmh:eq:central22log}
\end{align}
For direct reproduction of all $N=1$ cases, set $a=1/2+u$. The needed
one-factor Taylor series are
\begin{align}
 J_1(1/2+u,s)&=\pi\sec(\pi u)\biggl[
 (2-2^s)\zeta(s)
 -\sum_{k\ge1}\frac{(s)_k}{k!}(2^{s+k}-1)\zeta(s+k)u^k
 \biggr],\quad s\ge2,\label{bmh:eq:central-J}\\
 J_1(1/2+u,1)&=\pi\sec(\pi u)\biggl[
 -2\log2-\sum_{k\ge1}(2^{k+1}-1)\zeta(k+1)u^k
 \biggr].\label{bmh:eq:central-J1}
\end{align}
Only finitely many coefficients are required in any fixed instance of
\eqref{bmh:eq:central}. The supplied exact evaluator constructs these
coefficients, rather than differentiating an expression containing
$\zeta(1)$.

\section{Stieltjes order jets and continuous primitive combinations}
\label{bmh:sec:stieltjes}

\subsection{A bilinear identity at every Stieltjes index}
Recall that $\Li_1^{[r]}(z)$ means
$\left.\partial_s^r\Li_s(z)\right|_{s=1}$.

\begin{theorem}[Bilinear polylogarithm--Stieltjes reciprocity]
\label{bmh:thm:Stieltjes}
For every integer $n\ge1$, every $r\ge0$, and $-1<\Rea a<1$,
\begin{align}
 &\int_0^\infty\frac{x^{a-1}\Li_1^{[r]}(-x)}{1+x}
      \left\{\Li_n(-x)+(-1)^n\Li_n(-1/x)\right\}\dd x\notag\\
 &\quad=(-1)^r P_n(\partial_a)
 \left\{\pi\csc(\pi a)\bigl[\gamma_r(1)-\gamma_r(1-a)\bigr]\right\}.
 \label{bmh:eq:Stieltjes}
\end{align}
The right side at $a=0$ is a combined removable germ. Arbitrary fixed
$a$-derivatives give further ordinary identities with powers of
$\log x$. At $r=0$, the right side is a finite polygamma expression.
\end{theorem}
\begin{proof}
By~\eqref{bmh:eq:inversion}, the expression in braces on the left is
$P_n(\log x)$. For complex $s$ near one, the one-factor master already
proved in the repository~\cite{MLD} gives
\[
 \int_0^\infty\frac{x^{a-1}\Li_s(-x)}{1+x}\dd x
 =\pi\csc(\pi a)\{\zeta(s)-\zeta(s;1-a)\}.
\]
Both sides are holomorphic for $-1<\Rea a<1$, with the cosecant pole
removed as a whole. The residues of the two zeta functions at $s=1$
are identical. Therefore~\eqref{bmh:eq:Stieltjes-convention} gives
\[
 \left.\partial_s^r\{\zeta(s)-\zeta(s;1-a)\}\right|_{s=1}
 =(-1)^r\{\gamma_r(1)-\gamma_r(1-a)\}.
\]
Apply $P_n(\partial_a)$ and use the dominated differentiation statement
of Proposition~\ref{bmh:prop:strip}. Since $\gamma_0=-\psi$, the last
assertion follows. This proof is a new bilinear assembly of the classical
inversion polynomial and the existing one-factor Stieltjes master; those
inputs are not claimed as new results.
\end{proof}

For example, its $n=1,r=1$ content can be written as
\begin{equation}\label{bmh:eq:Stieltjes-example}
 \int_0^\infty\frac{x^{a-1}\log x}{1+x}
             \left.\partial_s\Li_s(-x)\right|_{s=1}\dd x
 =\frac{\dd}{\dd a}\left\{
    \pi\csc(\pi a)[\gamma_1(1-a)-\gamma_1(1)]\right\}.
\end{equation}
The sign on $\gamma_1$ is fixed by the Laurent convention, not by an
informal substitution of $s=1$ into a pole. The identities involve
spectral derivatives of polylogarithms and generalized Stieltjes
constants; they do not assert that these derivatives lie in a finite
ordinary-polylogarithm algebra.

\subsection{Spectral Stieltjes jets of the connected coordinate}
For $A\ge2$ the bulk definition~\eqref{bmh:eq:Cbulk} is holomorphic in
$B$ in a neighborhood of one. Put
\begin{equation}\label{bmh:eq:scriptC}
 \mathfrak C_{A,r}(q)=
 \left.\partial_B^r\bmhC_{A,B}(q)\right|_{B=1}.
\end{equation}
Normal convergence gives the explicit logarithmic-double-sum formula
\begin{align}
 \mathfrak C_{A,r}(q)=(-1)^r\biggl[&
 \sum_{\nu>\mu\ge0}\frac{\log^r(\mu+1)}{(\nu+1)^A(\mu+1)}
 -\sum_{\nu>\mu\ge0}\frac{\log^r(\mu+q)}{(\nu+q)^A(\mu+q)}\notag\\
 &+\zeta(A;q)\{\gamma_r(q)-\gamma_r(1)\}\biggr].
 \label{bmh:eq:scriptC-series}
\end{align}
For $r=0$ the convention $\log^0(1)=1$ applies. The normal convergence
is immediate from an outer $\nu^{-A}$ factor times a fixed power of
$\log\nu$.

\begin{proposition}[Exact discrete primitive for Stieltjes sums]
\label{bmh:prop:Stieltjes-shift}
For $A\ge2$, $r\ge0$, and positive $q$,
\begin{equation}\label{bmh:eq:Stieltjes-shift}
 \mathfrak C_{A,r}(q+1)-\mathfrak C_{A,r}(q)
 =(-1)^{r+1}q^{-A}\{\gamma_r(q)-\gamma_r(1)\}.
\end{equation}
Consequently, for an integer $L\ge1$,
\begin{align}
 \sum_{j=0}^{L-1}\frac{\gamma_r(q+j)-\gamma_r(1)}{(q+j)^A}
 =(-1)^{r+1}\{
 \mathfrak C_{A,r}(q+L)-\mathfrak C_{A,r}(q)\}.
 \label{bmh:eq:Stieltjes-finite-sum}
\end{align}
At integer arguments this is compatible with the logarithmic generalized
harmonic sums
\begin{equation}\label{bmh:eq:Stieltjes-harmonic}
 \mathfrak C_{A,r}(L)=(-1)^r
 \sum_{1\le k<j\le L-1}\frac{\log^r k}{j^A k}.
\end{equation}
\end{proposition}
\begin{proof}
Differentiate~\eqref{bmh:eq:Cshift} $r$ times in $B$ at one.
The pole-cancelled spectral derivative is
$\partial_B^rD_B(q)|_{B=1}=(-1)^r(\gamma_r(q)-\gamma_r(1))$.
This proves~\eqref{bmh:eq:Stieltjes-shift}; summing it telescopes.
Alternatively differentiating the finite identity~\eqref{bmh:eq:Charmonic}
gives~\eqref{bmh:eq:Stieltjes-harmonic}. The elementary Hurwitz shift
underlying this law is classical. The connected coordinate supplies a
single compatible notation for its double-sum and Stieltjes forms.
\end{proof}

\subsection{A continuous primitive ladder}
A discrete primitive is not a continuous antiderivative. The latter can
also be obtained for a natural adjacent-index combination. Set
\begin{equation}\label{bmh:eq:U}
 U_k(q)=\begin{cases}\zeta(k;q),&k\ge2,\\-\psi(q),&k=1.
 \end{cases}
\end{equation}
Then $U_k'(q)=-k\zeta(k+1;q)$, including $k=1$.

\begin{proposition}[Anchored adjacent-index primitive]
\label{bmh:prop:continuous-primitive}
For $A\ge2$, $B\ge1$, and positive $q,q_0$,
\begin{align}
 &\int_{q_0}^{q}\left\{A\bmhC_{A+1,B}(t)+B\bmhC_{A,B+1}(t)\right\}\dd t
 \notag\\
 &\quad=\bmhC_{A,B}(q_0)-\bmhC_{A,B}(q)
 +\{A\zeta(A+1,B)+B\zeta(A,B+1)\}(q-q_0)\notag\\
 &\qquad\quad+
 \frac{B\zeta(B+1)}{A-1}
       \{U_{A-1}(q)-U_{A-1}(q_0)\}.
 \label{bmh:eq:continuous-primitive}
\end{align}
All terms are ordinarily finite, including the polygamma case $A=2$.
\end{proposition}
\begin{proof}
Differentiate~\eqref{bmh:eq:Cbulk}, use~\eqref{bmh:eq:Zq-derivative},
and regroup the two shifted connected coordinates. This gives
\begin{align}
 \bmhC_{A,B}'(q)={}&-A\bmhC_{A+1,B}(q)-B\bmhC_{A,B+1}(q)\notag\\
 &+A\zeta(A+1,B)+B\zeta(A,B+1)
       -B\zeta(B+1)\zeta(A;q).
 \label{bmh:eq:C-differential}
\end{align}
Integrate and use the derivative of $U_{A-1}$. The anchoring differences
in~\eqref{bmh:eq:continuous-primitive} fix every integration constant.
At $A=2$, the correct $U_1=-\psi$ is the finite part of the Hurwitz
zeta function at one; inserting an undefined $\zeta(1;q)$ would obscure
an otherwise finite antiderivative.
\end{proof}

\subsection{An independent unequal-scale elementary example}
The finite-alphabet theorem does not imply that every two-scale example
requires genuine depth two. For $z>0$, put $r=1-1/z$. Then
\begin{equation}\label{bmh:eq:multiscale-example}
 \boxed{\displaystyle
 \int_0^\infty\frac{\log(1+x)\log(1+zx)}{(1+x)^2}\dd x
 =\frac{\log z+\Li_2(1-1/z)}{1-1/z}.}
\end{equation}
At $z=1$ the removable value is $2$. To prove the formula, let
$u=(1+x)^{-1}$ and use
$\log(1+zx)=\log z-\log u+\log(1-ru)$.
The integral becomes
\[
 \log z+2-\int_0^1\log u\log(1-ru)\dd u.
\]
For $|r|<1$, termwise integration gives
\[
 \int_0^1\log u\log(1-ru)\dd u
 =\sum_{k\ge1}\frac{r^k}{k(k+1)^2}.
\]
The partial fraction identity
$1/(k(k+1)^2)=1/k-1/(k+1)-1/(k+1)^2$
sums this series and yields~\eqref{bmh:eq:multiscale-example}.
Both sides are real analytic for $r<1$, so continuation proves it for
all $z>0$. In particular, at $z=2$ the value is
\begin{equation}\label{bmh:eq:multiscale2}
 \frac{\pi^2}{6}+2\log2-\log^22.
\end{equation}
The example illustrates why a general upper-bound theorem must not be
rephrased as a minimal-depth claim.

\section{Audit, integration status, and further research}
\label{bmh:sec:audit}

\subsection{What is proved and what is unchanged}
The principal analytic results of this article are exact ordinary
integral identities. Their proofs establish the necessary convergence,
parameter differentiation, cancellation, and finite coefficient
extraction. The delivered finite tests check implementations of these
arguments; they are not a substitute for the arguments.

The sampled source material did not reveal a new false equation in the
one-factor theorem or the cited canonical depth results. Accordingly we
do not manufacture a mathematical erratum. The concrete repository
change is a \emph{status refinement} for the tenth research question in
\src{sections/06-audit-research.tex}: its positive-integer-order finite-
closure sector now has a constructive solution, while the arbitrary-
spectral-order and minimal-depth portions remain open.

The current Gaussian $S_6$ and $S_8$ identities are not proved or
refuted by this work. The prior exact $S_4$ result is not counted again
as a new result. No claim is made that the binary incoming archives have
all been read, that all new claims in those archives have been reconciled,
or that these proofs have been checked in a proof assistant.

\subsection{Normalization corrections to avoid during integration}
The following are explicit safeguards for the new continuation, not
accusations that the inspected source made each mistake.
\begin{enumerate}[leftmargin=2em,itemsep=5pt]
\item \textbf{Use the product strip.} For two factors the lower condition
is $\Rea a>-2$, not $\Rea a>-1$. The value at $a=-1$ is an ordinary
integral and has the exact formula~\eqref{bmh:eq:negative-jets}.
\item \textbf{Use connected leading-one coordinates.} The separate
objects $\zeta(1,B)$ and $Z_q(1,B)$ generally diverge. Equations
\eqref{bmh:eq:Cedge}--\eqref{bmh:eq:Ccorner} give convergent replacements,
derived by a combined analytic continuation.
\item \textbf{Keep the zero-order compact value correct.}
$\Li_0(-t/(1-t))=-t$. Using $-t/(1-t)$ breaks the kernel recurrence,
including its simplest beta base.
\item \textbf{Cancel poles before differentiating a specialization.}
Examples include $a=0,-1$, $s=1$, and the one-factor product
$s\zeta(s+1)$ at $s=0$. The last removable value is $1$, not $0$.
\item \textbf{Separate algebra membership from additive depth.}
Even-zeta factors times double-zeta values belong to the algebra generated
by doubles, but can have additive depth three. The word-weight filtration
likewise supplies no theorem of numerical independence.
\item \textbf{Do not call a discrete primitive an antiderivative.}
Equation~\eqref{bmh:eq:Stieltjes-finite-sum} is a shift identity;
\eqref{bmh:eq:continuous-primitive} and~\eqref{bmh:eq:primitive} are
continuous anchored primitives with explicit constants.
\end{enumerate}

\subsection{Proposed research questions}
The questions below seek exact identities or structural classifications,
not numerical lower bounds inferred from failed searches.
\begin{enumerate}[label=\textbf{Q\arabic*.},leftmargin=2.8em,itemsep=7pt]
\item \textbf{A genuinely two-scale double-Hurwitz formula.}
For independent $z,w$, replace the unit-scale cone decomposition by a
shifted-cone calculation and identify a minimal convenient family of
multiple Lerch coordinates. The answer should specialize to
\eqref{bmh:eq:master} when $z=w=1$ and to
\eqref{bmh:eq:multiscale-example} at logarithmic order.
\item \textbf{Three or more equal-scale factors.}
A product of three Fermi integrals splits into six order cones. Determine
a finite connected triple-Hurwitz system that makes all leading-one
faces convergent. The two-factor coordinates suggest an inclusion--
exclusion pattern, but its exact signs and edge terms require proof.
\item \textbf{Depth of the weighted MZV family.}
The specific weighted sums in~\eqref{bmh:eq:weighted} reduce to double
values at every weight. Find a short double-shuffle proof that does not
pass through the integral, and determine which further coefficient
symmetries force a single-zeta reduction. Any arithmetic minimal-depth
claim requires additional input and must be stated separately.
\item \textbf{The missing central parity sector.}
When $W+h$ is even, inversion removes the desired central moment from
\eqref{bmh:eq:reciprocity}. Determine an independent functional relation
for this complementary sector. Further differentiation of the same
projection cannot recover a coefficient that it annihilates.
\item \textbf{A complete continuous primitive system.}
Equation~\eqref{bmh:eq:continuous-primitive} integrates a specified
adjacent-index combination. Construct a basis in which individual
connected coordinates have explicit primitives, with all integration
constants and leading-one specializations controlled.
\item \textbf{Spectral derivatives of the double-Hurwitz master.}
The present cone expansion uses integer $m,n$. Find a representation at
complex $m,n$ that is differentiable in the spectral variables and
specializes to the finite binomial formula. This is a precise route toward
products containing $\Li_s^{[r]}$ without assuming finite ordinary-MPL
closure at noninteger order.
\item \textbf{Rational common shifts and cyclotomic coordinates.}
At $q=p/d>0$, express the connected double-Hurwitz values and their
Stieltjes spectral jets through residue-class filters. Quantify the finite
alphabet and the exact cancellation at leading-one coordinates before
reducing Dirichlet-character combinations.
\item \textbf{Higher negative resonances for more factors.}
An $r$-factor product has ordinary integer values at
$a=1-r,\ldots,-1$. Determine whether the connected higher-depth system
has a shift law as simple as~\eqref{bmh:eq:Cshift}, and derive every
logarithmic moment at those resonances.
\item \textbf{A proof-assistant boundary.}
The closed shuffle multiplicities, Laurent decomposition, finite
recurrences, and coefficient convolutions are natural first targets for
formalization. The remaining analytic obligations are explicit: dominated
holomorphy, the Fermi integral, the combined cone limit, and endpoint
regularization. No finite test count alone discharges these obligations.
\end{enumerate}

\subsection{Integration recommendation}
Place this delivery as a continuation of the Mellin--Lerch report under
\src{continuations/bilinear-mellin-hurwitz/}, or under a report path chosen
by the repository intake process. Preserve the self-contained article and
its checks. The canonical integration chapter can adopt the master,
resonance, and kernel sections after analytic review; the differentiation
chapter can adopt the Stieltjes identities, and the depth chapter can
adopt the weighted MZV theorem. Every theorem and equation label here has
the prefix \src{bmh:}. The accompanying integration note supplies the
precise scope wording for the research-question update.

\appendix
\section{Reproducible finite verification and numerical diagnostics}
\label{bmh:sec:verification}

\subsection{What was executed}
The package contains two independent implementation layers. The exact
layer uses Python integers and SymPy rational arithmetic, with an
unevaluated symbol $Z(A,B)$ for convergent double zeta values. The
numerical layer uses mpmath quadrature of the original Mellin integrals
and an independently implemented finite-sum plus Euler--Maclaurin
representation of double Hurwitz zeta values. It does not evaluate an
integral by merely calling the symbolic result being tested.

The exact run passed \textbf{8,849 checks}. Its recorded categories are
\begin{center}
\begin{tabular}{@{}lr@{}}
\toprule
Exact category & Checks\\
\midrule
Enumerated shuffle multiplicities through word length ten & 8,194\\
Cone formula versus compact coefficients & 91\\
Compact-formula symmetry & 91\\
Laurent decomposition, endpoint, transport, residue identities & 216\\
Gamma product versus independent Bell recurrence & 90\\
Negative resonance versus independent kernel recursion & 48\\
Inversion derivative and parity laws & 30\\
Central-moment symmetry & 85\\
Printed low-weight examples & 3\\
Deliberately corrupted coefficient rejected & 1\\
\midrule
Total & 8,849\\
\bottomrule
\end{tabular}
\end{center}
The shuffle test enumerates subsets of word positions independently of
the binomial formula; it does not merely evaluate that formula twice.
The negative-resonance test compares the common-shift Taylor law with the
separately implemented integration-by-parts recurrence. The corruption
control adds a nonzero formal double-zeta coefficient and verifies that
exact cancellation fails.

The numerical run passed \textbf{67 comparisons at 45 decimal working
digits}. The largest scaled discrepancy was less than
$2.3\times10^{-44}$, against an acceptance threshold $10^{-33}$.
The families tested were the noninteger master (nine cases), central
$a=0$ logarithmic moments (twelve), negative-resonance moments (five),
integer-kernel recurrence values (ten), central parity values (seven),
unequal-scale logarithmic examples (five), connected shifts and stuffles
(ten), independent double-zeta tail refinements (four), and beta moments
(five). The Stieltjes spectral-derivative and continuous-primitive
extensions are proved analytically in Section~\ref{bmh:sec:stieltjes};
the 67 numerical checks are not advertised as direct quadrature tests
of those extensions.

Representative raw values, printed here only as diagnostics, are
\begin{align*}
 I_{1,1}&\approx2.40411380631918857079947632302,\\
 I_{2,1}&\approx4.59987374327233731394301571029,\\
 I_{3,1}&\approx7.68040700358583071301959561699,\\
 I_{2,2}&\approx9.98307291147592432545107273885.
\end{align*}
The full run records, not these rounded lines, are the replay evidence.
They are floating-point diagnostics, not rational interval certificates.
An accepted tiny residual is not used anywhere as a proof of equality.

\subsection{The independent double-zeta evaluator}
For $A\ge2,B\ge2$, write
\[
 Z_q(A,B)=\zeta(A;q)\zeta(B;q)
       -\sum_{n\ge0}(n+q)^{-A}\zeta(B;n+q).
\]
After a finite sum through $n=M-1$, the tail is evaluated with
\begin{align*}
 &\frac{\zeta(A+B-1;M+q)}{B-1}
 +\frac12\zeta(A+B;M+q)\\
 &\qquad+\sum_{k=1}^{K}\frac{B_{2k}}{(2k)!}(B)_{2k-1}
          \zeta(A+B+2k-1;M+q).
\end{align*}
For $B=1$, use the finite partial sum of
$(n+q)^{-A}(\psi(n+q)-\psi(q))$ and the tail
\begin{align*}
 &-\zeta^{[1]}(A;M+q)-\psi(q)\zeta(A;M+q)
   -\frac12\zeta(A+1;M+q)\\
 &\qquad-\sum_{k=1}^{K}\frac{B_{2k}}{2k}
                      \zeta(A+2k;M+q).
\end{align*}
Here $\zeta^{[1]}$ differentiates the spectral variable. The default
parameters are $M=64,K=24$; a second calculation uses $M=96,K=32$.
These are truncated asymptotic tails checked by refinement, not asserted
convergent infinite expansions or certified error bounds.

The direct quadrature uses $u=\log x$. Polylogarithm inversion keeps
integer-order library evaluations inside the unit disk. At order one,
$-\log(1+e^u)$ is computed with a stable \src{log1p} form; a naive
$\log(1+\varepsilon)$ can lose the very small lower-tail contribution at
negative Mellin parameters. This implementation detail does not alter
the mathematical integral.

\subsection{Commands and delivery contents}
From the package root, with Python, SymPy, mpmath, and pdfLaTeX installed:
\begin{verbatim}
python verification/verify_exact.py
python verification/verify_numerics.py
python verification/export_tables.py
python build.py
\end{verbatim}
The tested Python version is 3.13.5, SymPy is 1.14.0, and the numerical
record identifies the mpmath version. The supplied requirements pin the
Python packages. Alternative output paths can be passed to the two
verification scripts with \src{--output}; this permits replay on a copy
without overwriting delivered evidence. The scripts fail with a nonzero
exit status when an assertion fails.

The article is a self-contained source split into section files.
\src{data/exact_results.json} and \src{data/numerical_results.json}
record the completed runs; \src{data/formula_table.json} stores exact
example expressions; \src{AUDIT.md}, \src{PROVENANCE.json}, and
\src{integration/INTEGRATION.md} record status and source scope.
The build script checks the LaTeX process and runs enough passes to
resolve cross-references. The delivered PDF was also rendered and
visually inspected. No proof-assistant certification or external peer
review is claimed.

\begin{thebibliography}{9}
\bibitem{Repo}
ProveIt Contributors, \emph{Polylogarithms and their Arithmetic Bridges:
Special Values, Zeta Jets, Lattices and Experimental Discovery}, collective
manuscript, October 2026. Repository:
\url{https://github.com/VladimirReshetnikov/ProveIt}.
Source directory: \url{Analysis/Polylogarithms/docs/manuscript/}.
The exact inspected scope is recorded in the delivery provenance.

\bibitem{MLD}
ProveIt research continuation, \emph{Mellin--Lerch Transforms and Dilation
Identities for Stieltjes Correlations: Polylogarithms, Generalized Harmonic
Sums, and Exact Finite Parts}, 10 October 2026, prepared with OpenAI
assistance. Repository report under
\url{Analysis/Polylogarithms/docs/reports/stieltjes-correlation-calculus/continuations/mellin-lerch-dilation/}.
In particular Sections 2--3 and the tenth research question of
\src{sections/06-audit-research.tex}. The source report records its own
original pin \src{fc4d3bf80534ad7c901d3b8c9e71baf2df0064ed}.

\bibitem{Panzer}
E. Panzer, \emph{Algorithms for the symbolic integration of hyperlogarithms
with applications to Feynman integrals}, Computer Physics Communications
\textbf{188} (2015), 148--166. Preprint arXiv:1403.3385.
\url{https://arxiv.org/abs/1403.3385}.
The rational-hyperlogarithm integration mechanism is in Section 2.2;
its use here is a specialized application, not a new general algorithm.

\bibitem{DLMFhurwitz}
NIST Digital Library of Mathematical Functions, \S25.11, \emph{Hurwitz
Zeta Function}. \url{https://dlmf.nist.gov/25.11}.
Used for standard conventions, parameter derivatives, and relations to
polygamma functions. Accessed 10 October 2026.

\bibitem{DLMFpoly}
NIST Digital Library of Mathematical Functions, \S25.12, \emph{Polylogarithms}.
\url{https://dlmf.nist.gov/25.12}. Used for the classical series and integral
representations and branch conventions. Accessed 10 October 2026.

\bibitem{DLMFbeta}
NIST Digital Library of Mathematical Functions, \S5.12, \emph{Beta Function},
and \S5.5, \emph{Functional Relations}.
\url{https://dlmf.nist.gov/5.12}; \url{https://dlmf.nist.gov/5.5}.
Used for the beta integral, Gamma quotient, and reflection normalization.
Accessed 10 October 2026.
\end{thebibliography}

\end{document}
